\chapter{Introducción}

Este trabajo aborda la optimización puntual de la planificación académica de la Facultad de Ingeniería mediante técnicas de grafos implementadas en \textbf{TypeScript}. El problema tratado es la asignación de salones y franjas horarias a bloques académicos evitando conflictos entre estudiantes, profesores y tipos de espacio.

La solución se organiza en capas: carga y normalización de datos, construcción del grafo de conflictos, algoritmos de resolución (coloreo y caminos mínimos), generación del horario final, exposición por API y visualización en frontend. Todas las decisiones de diseño responden a criterios de reproducibilidad y trazabilidad del código.

Objetivo específico: producir una asignación válida de salones para los bloques académicos del dataset, comparar al menos tres heurísticas de coloreo y generar horarios personalizados por estudiante usando técnicas de caminos mínimos cuando sea necesario.

